% Part: computability
% Chapter: computability-theory
% Section: fixed-point-thm

\documentclass[../../../include/open-logic-section]{subfiles}

\begin{document}

\olfileid{cmp}{thy}{fix}
\olsection{स्थिरबिंदू प्रमेय}

थांबण्याच्या समस्येचा पुन्हा विचार करू. तात्पुरत्या चिन्हांकनासाठी,
$\tuple{x, y}$ ऐवजी $\gn{\cfind{x}(y)}$ असे लिहू;
याला $\cfind{x}(y)$ या मूल्याचे ``नाव'' दर्शवणारे चिन्ह समजा.
या चिन्हांकनाने थांबण्याची समस्या अनिर्णेय असल्याचा आपला एक
पुरावा पुन्हा मांडता येतो.

प्रश्न: पुढील गुणधर्म असलेले संगणनक्षम फलन $h$ आहे का?
प्रत्येक $x$ आणि $y$ साठी,
\[
h(\gn{\cfind{x}(y)}) =
\begin{cases}
1 & \text{जर $\cfind{x}(y) \fdefined$ असेल तर} \\
0 & \text{अन्यथा.}
\end{cases}
\]

उत्तर: नाही. अन्यथा पुढील आंशिक फलन
\[
g(x) \simeq
\begin{cases}
0 & \text{जर $h(\gn{\cfind{x}(x)}) = 0$ असेल तर} \\
\text{अपरिभाषित} & \text{अन्यथा}
\end{cases}
\]
संगणनक्षम असते आणि त्याला काही निर्देशांक~$e$ असता.
पण मग
\[
\cfind{e}(e) \simeq
\begin{cases}
0 & \text{जर $h(\gn{\cfind{e}(e)}) = 0$ असेल तर} \\
\text{अपरिभाषित} & \text{अन्यथा,}
\end{cases}
\]
असे मिळते. त्यानुसार $\cfind{e}(e)$ परिभाषित असेल तेव्हाच
ते अपरिभाषित असेल; हा विरोधाभास आहे.

आता $\cfind{e}$ असलेल्या समीकरणाकडे पाहा. त्यात एका
अर्थाने स्वसंदर्भ आहे: $\cfind{e}(e)$ चे मूल्य विशिष्ट
पद्धतीने $\gn{\cfind{e}(e)}$ वर अवलंबून राहील अशी
रचना केली आहे. स्थिरबिंदू प्रमेय सांगते की हे सर्वसाधारणपणे
करता {\em येते}—फक्त विरोधाभास सिद्ध करण्यासाठीच नाही.

\olref{lem:fixed-equiv} मध्ये स्थिरबिंदू प्रमेय मांडण्याच्या
दोन समतुल्य पद्धती दिल्या आहेत. तार्किकदृष्ट्या दोन्ही
विधाने खरी असल्याने ती समतुल्य आहेत असे म्हणता येते;
पण आपला खरा अर्थ असा आहे की प्रत्येक विधानावरून दुसरे
सहज निष्पन्न होते. म्हणून ती एकाच प्रमेयाची पर्यायी विधाने
मानता येतात.

\begin{lem}
\ollabel{lem:fixed-equiv}
पुढील विधाने समतुल्य आहेत:
\begin{enumerate}
\item प्रत्येक आंशिक संगणनक्षम फलन $g(x,y)$ साठी असा
  निर्देशांक~$e$ आहे की प्रत्येक~$y$ साठी
\[
\cfind{e}(y) \simeq g(e,y).
\]
\item प्रत्येक संगणनक्षम फलन~$f(x)$ साठी असा निर्देशांक~$e$
  आहे की प्रत्येक~$y$ साठी
\[
\cfind{e}(y) \simeq \cfind{f(e)}(y).
\]
\end{enumerate}
\end{lem}

\begin{proof}
$(1) \Rightarrow (2)$: $f$ दिल्यावर $g(x,y) \simeq
\fn{Un}(f(x),y)$ असे $g$ व्याख्यायित करा. (1) वापरून
प्रत्येक~$y$ साठी पुढील गुणधर्म असलेला निर्देशांक~$e$
मिळवा:
\begin{align*}
\cfind{e}(y) & = \fn{Un}(f(e),y) \\
& = \cfind{f(e)}(y).
\end{align*}

$(2) \Rightarrow (1)$: $g$ दिल्यावर $s$-$m$-$n$
प्रमेय वापरून प्रत्येक $x$ आणि~$y$ साठी
$\cfind{f(x)}(y) \simeq g(x,y)$ होईल असे $f$ मिळवा.
(2) वापरून असा निर्देशांक~$e$ मिळवा की
\begin{align*}
\cfind{e}(y) & = \cfind{f(e)}(y) \\
& = g(e,y).
\end{align*}
याने सिद्धता पूर्ण झाली.
\end{proof}

\begin{explain}
विधान (1) खरे आहे (आणि म्हणून (2) देखील) हे दाखवण्यापूर्वी,
ते किती विलक्षण आहे याचा विचार करा. $e$ हा संगणक
कार्यक्रम आहे असे समजा; विधान (1) सांगते की कोणतेही
आंशिक संगणनक्षम $g(x,y)$ दिले असता, $g_e(y) \simeq g(e,y)$
संगणित करणारा संगणक कार्यक्रम $e$ शोधता येतो. म्हणजे
स्वतःचाच संदर्भ असलेले फलन संगणित करणारा कार्यक्रम शोधता येतो.
\end{explain}

\begin{thm}
\olref{lem:fixed-equiv} मधील दोन्ही विधाने खरी आहेत.
विशेषतः, प्रत्येक आंशिक संगणनक्षम फलन $g(x,y)$ साठी
असा निर्देशांक~$e$ आहे की प्रत्येक~$y$ साठी
\[
\cfind{e}(y) \simeq g(e,y).
\]
\end{thm}

\begin{proof}
आवश्यक घटक वरच्या थांबण्याच्या समस्येवरील चर्चेत अप्रत्यक्षपणे
आहेतच. $\fn{diag}(x)$ हे असे संगणनक्षम फलन असू द्या की
ते प्रत्येक $x$ साठी $f_x(y) \simeq \cfind{x}(x,y)$
या फलनाचा निर्देशांक परत करते, म्हणजे
\[
\cfind{\fn{diag}(x)}(y) \simeq \cfind{x}(x,y).
\]
$\fn{diag}$ कडे द्विस्थानी फलनाचा कार्यक्रम एकस्थानी
फलनाच्या कार्यक्रमात रूपांतरित करणारे फलन म्हणून पाहा;
मूळ कार्यक्रमाचा पहिला युक्तिवाद त्याच कार्यक्रमावर निश्चित
करून हे फलन मिळते. $\fn{diag}$ ची औपचारिक व्याख्या अशी
करता येते: आधी $s$ ची व्याख्या
\[
s(x,y) \simeq \fn{Un}^2(x,x,y),
\]
अशी करा. येथे $\fn{Un}^2$ हे आंशिक
संगणनक्षम द्विस्थानी फलनांसाठी सार्वत्रिक असलेले त्रिस्थानी
फलन आहे. मग $s$-$m$-$n$ प्रमेयानुसार पुढील अट पूर्ण करणारे
आदिम पुनरावर्ती फलन $\fn{diag}$ मिळते:
\[
\cfind{\fn{diag}(x)}(y) \simeq s(x,y).
\]

आता $l$ फलनाची व्याख्या अशी करा:
\[
l(x,y) \simeq g(\fn{diag}(x),y).
\]
आणि $\gn{l}$ हा $l$ चा निर्देशांक असू द्या. शेवटी
$e = \fn{diag}(\gn{l})$ असू द्या. मग प्रत्येक $y$ साठी
\begin{align*}
\cfind{e}(y) & \simeq \cfind{\fn{diag}(\gn{l})}(y) \\
& \simeq \cfind{\gn{l}}(\gn{l}, y) \\
& \simeq l(\gn{l}, y) \\
& \simeq g(\fn{diag}(\gn{l}),y) \\
& \simeq g(e, y),
\end{align*}
हवे तसे.
\end{proof}

\begin{explain}
येथे नेमके काय घडते? स्वतःचाच मजकूर छापणारा संगणक
कार्यक्रम लिहायचा आहे असे समजा. यासाठी तुमची प्रोग्रामिंग
भाषा चिन्हमालांवरील समृद्ध व काहीशी विचित्र फलनांचा संच
पुरवते असेही समजा. विशेषतः, तिच्यात पुढीलप्रमाणे काम करणारे
फलन $\fn{diag}$ आहे असे समजा: आदान चिन्हमाला~$s$
दिली असता $\fn{diag}$ $s$ मधील `x' या चिन्हाच्या प्रत्येक
घटनेचा शोध घेते आणि त्या जागी मूळ चिन्हमालेची अवतरणचिन्हांत
घेतलेली आवृत्ती बसवते. उदाहरणार्थ, पुढील चिन्हमाला
\begin{quote}
\begin{verbatim}
hello x world
\end{verbatim}
\end{quote}
आदान म्हणून दिल्यास हे फलन पुढील चिन्हमाला परत करते:
\begin{quote}
\begin{verbatim}
hello 'hello x world' world
\end{verbatim}
\end{quote}
मग अपेक्षित कार्यक्रम लिहिणे सोपे आहे; पुढील कार्यक्रम
हे काम करतो हे तपासा:
\begin{quote}
\begin{verbatim}
print(diag('print(diag(x))'))
\end{verbatim}
\end{quote}
C++ आणि Java सारख्या अधिक प्रचलित प्रोग्रामिंग भाषांतही
हीच कल्पना (अधिक गुंतागुंतीच्या अंमलबजावणीसह) काम करते.

आता आपण स्थिरबिंदू प्रमेयाच्या सिद्धतेपासून काही पावलेच
दूर आहोत. $\fn{print}(x,y)$ हे print फलनाचे एक रूप
चिन्हमाला $x$ आणि संख्यात्मक युक्तिवाद $y$ घेते व
$x$ ही चिन्हमाला $y$ वेळा छापते असे समजा. मग पुढील ``कार्यक्रम''
\begin{quote}
\begin{verbatim}
getinput(y);
print(diag('getinput(y); print(diag(x), y)'), y)
\end{verbatim}
\end{quote}
आदान $y$ मिळाल्यावर स्वतःचा मजकूर $y$ वेळा छापतो.
$\fn{getinput}$---$\fn{print}$---$\fn{diag}$ या सांगाड्याच्या
जागी कोणतेही फलन $g(x,y)$ घेतल्यास
\begin{quote}
\begin{verbatim}
g(diag('g(diag(x), y)'), y)
\end{verbatim}
\end{quote}
मिळते. हा कार्यक्रम आदान $y$ वर स्वतःच्या कार्यक्रमाला
आणि $y$ ला $g$ चे युक्तिवाद म्हणून देतो. ``अवतरणचिन्हांत
घेणे'' म्हणजे ``निर्देशांक वापरणे'' असे मानल्यास वरील
सिद्धता मिळते.

सध्या हा पुरावा औपचारिक युक्ती किंवा जादू वाटला तरी
हरकत नाही. पण वरील युक्तिवादाचे तपशील पुन्हा उभारता
आले पाहिजेत. अपूर्णतेची प्रमेये (आणि संबंधित ``स्थिरबिंदू
प्रमेय'') सिद्ध करताना ते का काम करते हे समजण्याच्या
इतर पद्धतींवर चर्चा करू.
\end{explain}

\begin{tagblock}{lambda}
\begin{digress}
याच कल्पनेने ``स्थिरबिंदू'' संयोजक मिळवता येतो. $g$
हे लॅम्डा पद दिले आहे आणि दुसरे पद $k$ हवे आहे असे
समजा; $k$ हे $gk$ शी $\beta$-सममूल्य असावे. आपल्या चिन्हांकन
परंपरेनुसार पुढील पदे व्याख्यायित करा:
\[
\fn{diag}(x) = xx
\]
आणि
\[
l(x) = g(\fn{diag}(x))
\]
म्हणजे $l$ हे $\lambd[x][g(xx)]$ पद आहे. $k$ हे $ll$
पद असू द्या. मग
\begin{align*}
k & = (\lambd[x][g(xx)])(\lambd[x][g(xx)]) \\
& \red  g((\lambd[x][g(xx)])(\lambd[x][g(xx)])) \\
& = gk.
\end{align*}
आता
\[
Y = \lambd[g][((\lambd[x][g(xx)])(\lambd[x][g(xx)]))]
\]
असे घेतल्यास $Yg$ आणि $g(Yg)$ ही दोन्ही पदे एका
सामायिक पदापर्यंत न्यूनीत होतात; म्हणून $Yg \equiv_\beta
g(Yg)$. याला ``करीचा संयोजक'' म्हणतात. त्याऐवजी
\[
Y = (\lambd[xg][g(xxg)])(\lambd[xg][g(xxg)])
\]
घेतल्यास $Yg$ थेट $g(Yg)$ पर्यंत न्यूनीत होते;
हे अधिक मजबूत विधान आहे. $Y$ च्या या रूपाला
``ट्यूरिंगचा संयोजक'' म्हणतात.
\end{digress}
\end{tagblock}

\end{document}
